Besides the extension of the Script Assembler's capabilities and the creation of a new EBNF language that makes the software more extendable, another main focus of the project is the creation of a mockup of the TestRail database. This mockup is not only made for testing purposes, but allows the user to work independently of the TestRail instance. The mockup is designed to be as close to the real TestRail database as possible, allowing a seamless switch between the two instances. This means that all functionalities that are available in the real TestRail database are also available in the new mockup.

This section aims to give an overview of the different types of databases, the structure and functionalities of the TestRail database specifically, the resulting design and implementation of the mockup database, and the way the database switch is finally implemented.

\section{Types of Databases}

Depending on the requirements of a given project, different types of databases are suitable for different use cases. It is important to analyze what qualities are needed and which qualities are negligible to find the best fitting solution. 

The most important qualities of a database type can be expressed in the following categories:

\begin{itemize}
    \item \textbf{Data Model:} \\
    The data model of a certain database type describes how the data is structured and stored. This can be in the form of tables, documents, graphs or other structures. The data model has a notable impact on aspects like performance and scalability.

    \item \textbf{Storage Architecture:} \\
    The storage architecture of a database type not only describes how the data is stored, but also how it may be accessed, managed and distributed. This can be in the form of a single server, a cluster of servers or even a distributed network of servers across the globe.

    \item \textbf{Query Model:} \\
    The query model of a database type describes the means of accessing the data. This can occur in the form of a specifically designed query language, an API or other methods. This has an undeniable impact on the ease of use of the database and how the user writes and optimizes queries.
\end{itemize}

Based on these qualities, various databases types have been developed, each with their own unique advantages and disadvantages. Some of the most common types of databases include but are not limited to relational databases, object-oriented databases and NoSQL databases, which can be further categorized by the data model they use, such as document stores, key-value stores, column-family stores and graph databases. (cf.~\cite{DigitalOceanDatabaseTypes})

\section{Relational Databases}

\section{The TestRail Database}

\section{The Mockup Database}

\section{Switching Between Databases}